\section{Composizione di Funzioni e Callback (Lez. 10)}

\subsection{Comporre le Funzioni}
In programmazione, spesso si usa l'output di una funzione come input per un'altra. Questo si chiama \textbf{composizione}.

\begin{concept}{Matrioska di Funzioni}
In programmazione, spesso si usa l'output di una funzione come input per un'altra. Questo si chiama \textbf{composizione}.
$f(g(x))$ in matematica.
\end{concept}
\begin{jsbox}
function raddoppia(x) { return x * 2; }
function quadrato(x) { return x * x; }

// Composizione manuale: prima il quadrato, poi raddoppio
let risultato = raddoppia(quadrato(3)); 
// Passo 1: 3 * 3 = 9
// Passo 2: 9 * 2 = 18
console.log(risultato); // 18
\end{jsbox}

\subsection{Funzioni che ritornano Funzioni (Function Factory)}
\begin{concept}{Fabbrica di Funzioni}
Possiamo scrivere funzioni che \textit{creano} nuove funzioni personalizzate, "congelando" alcuni parametri (Closure).
\end{concept}
\begin{jsbox}
function creaPolinomio(a, b, c) {
    // Ritorna una NUOVA funzione che accetta solo x
    return function(x) {
        return (a * x * x) + (b * x) + c;
    }
}

// Creo una funzione specifica: 2x^2 + x + 0
let miaEquazione = creaPolinomio(2, 1, 0);

// Ora la uso solo con x
console.log(miaEquazione(3)); // (2*9) + 3 + 0 = 21
console.log(miaEquazione(10)); // (2*100) + 10 + 0 = 210
\end{jsbox}
